\documentclass[11pt,a4paper]{article}

% Typeface: XeLaTeX with Latin Modern Roman (fontspec default), British English
\usepackage{fontspec}
\usepackage{polyglossia}
\setdefaultlanguage[variant=british]{english}

% Page: the measures of the note template (A4, 11 pt, 4.5 cm side margins, header and page count)
\usepackage[a4paper,left=4.5cm,right=4.5cm,top=4.5cm,bottom=4.1cm,headheight=14pt,headsep=0.55cm,
  footskip=0.55cm]{geometry}
\linespread{0.946}
\usepackage{fancyhdr}
\usepackage{lastpage}
\pagestyle{fancy}
\fancyhf{}
\fancyhead[L]{\small\textbf{PROJECT NOTE} $|$ bitcoin-macro-blueml}
\fancyhead[R]{\small\textbf{Compiled on:} \notedate}
\fancyfoot[C]{\small\thepage\,/\,\pageref{LastPage}}
\renewcommand{\headrulewidth}{0.4pt}

\usepackage{microtype}
\usepackage{amsmath,amssymb}
\usepackage{booktabs}
\usepackage{graphicx}
\usepackage{enumitem}
\setlist{noitemsep, topsep=3pt, leftmargin=*}
\usepackage{xcolor}
\definecolor{secblue}{HTML}{1F3A5F}
\definecolor{boxback}{HTML}{EEF2F6}
\definecolor{mark}{HTML}{9B2226}
\definecolor{hl}{HTML}{FFF36B}
% highlighter: each word in a yellow box, line breaks allowed between words
\makeatletter
\def\hl#1{{\setlength{\fboxsep}{0.4pt}\hl@word#1 \@nil}}
\def\hl@word#1 #2\@nil{\colorbox{hl}{\strut #1}%
  \if\relax\detokenize{#2}\relax\else\colorbox{hl}{\strut\ }\allowbreak\hl@word#2\@nil\fi}
\makeatother
\usepackage{titlesec}
\titleformat{\section}{\color{secblue}\bfseries\normalsize}{\thesection.}{0.4em}{}
\titlespacing*{\section}{0pt}{10pt}{4pt}
\usepackage[most]{tcolorbox}
\newtcolorbox{notebox}[1]{colback=boxback, colframe=secblue, boxrule=0.5pt, arc=0pt, left=4pt, right=4pt,
  top=3pt, bottom=3pt, fonttitle=\bfseries\color{secblue}, coltitle=secblue, colbacktitle=boxback,
  title={#1}, toptitle=1pt, bottomtitle=0pt}
\usepackage[font=small,labelfont=bf,labelsep=period,skip=4pt]{caption}
\usepackage[bottom]{footmisc}
\setlength{\parindent}{0pt}
\setlength{\parskip}{5pt}
\setlength{\textfloatsep}{10pt}
\emergencystretch=3em
\newcommand{\name}[1]{\textcolor{mark}{\textbf{#1}}}
\newcommand{\widefig}[1]{\includegraphics[width=\textwidth]{#1}}

\newcommand{\notedate}{\number\day\space\ifcase\month\or January\or February\or March\or April\or May\or June\or
  July\or August\or September\or October\or November\or December\fi\space\number\year}

\renewcommand{\topfraction}{0.85}\renewcommand{\bottomfraction}{0.7}\renewcommand{\textfraction}{0.1}\renewcommand{\floatpagefraction}{0.8}
\newcommand{\Months}{127}
\newcommand{\BtcFirst}{218.5}
\newcommand{\BtcLast}{115,765}
\newcommand{\SpxFirst}{1,995}
\newcommand{\SpxLast}{6,339}
\newcommand{\BtcGrowth}{52,882}
\newcommand{\SpxGrowth}{218}
\newcommand{\CorrLevels}{0.94}
\newcommand{\CorrReturns}{0.36}
\newcommand{\LevelAdfPmin}{0.13}
\newcommand{\CpiDiffAdfP}{0.37}
\newcommand{\CpiDiffPpP}{$<$0.001}
\newcommand{\InflAdfP}{0.32}
\newcommand{\InflAdfLags}{10}
\newcommand{\InflPpP}{$<$0.001}
\newcommand{\RetAdfPmax}{$<$0.001}
\newcommand{\GrangerPmin}{0.19}
\newcommand{\RsqOne}{0.223}
\newcommand{\RsqTwo}{0.155}
\newcommand{\Nobs}{126}
\newcommand{\SpxOne}{16.13}
\newcommand{\SpxTwo}{1.67}
\newcommand{\SpxTwoSeCl}{0.39}
\newcommand{\SpxOneSeCl}{2.97}
\newcommand{\SpxTwoSeHc}{0.36}
\newcommand{\SpxOneSeHc}{3.75}
\newcommand{\SpxTwoSeHac}{0.33}
\newcommand{\SpxOneSeHac}{3.71}
\newcommand{\FedOneP}{0.128}
\newcommand{\FedOnePHc}{0.042}
\newcommand{\ConstTwo}{4.77}
\newcommand{\ConstTwoP}{0.033}
\newcommand{\OthersTwoPmin}{0.27}
\newcommand{\ResetOneP}{0.569}
\newcommand{\ResetOne}{0.57}
\newcommand{\VifOne}{1.73}
\newcommand{\CondOne}{2.31}
\newcommand{\BpOneP}{0.039}
\newcommand{\BpOne}{11.74}
\newcommand{\WhiteOneP}{0.013}
\newcommand{\WhiteOne}{36.70}
\newcommand{\BgTwoOneP}{0.008}
\newcommand{\BgTwoOne}{9.60}
\newcommand{\BgSixOneP}{0.009}
\newcommand{\BgSixOne}{17.14}
\newcommand{\BgTwelveOneP}{0.102}
\newcommand{\BgTwelveOne}{18.46}
\newcommand{\DwOne}{1.67}
\newcommand{\JbOneP}{$<$0.001}
\newcommand{\JbOne}{107.9}
\newcommand{\ResetTwoP}{0.356}
\newcommand{\ResetTwo}{1.04}
\newcommand{\VifTwo}{1.67}
\newcommand{\CondTwo}{2.23}
\newcommand{\BpTwoP}{0.844}
\newcommand{\BpTwo}{2.04}
\newcommand{\WhiteTwoP}{0.837}
\newcommand{\WhiteTwo}{13.87}
\newcommand{\BgTwoTwoP}{0.421}
\newcommand{\BgTwoTwo}{1.73}
\newcommand{\BgSixTwoP}{0.940}
\newcommand{\BgSixTwo}{1.76}
\newcommand{\BgTwelveTwoP}{0.648}
\newcommand{\BgTwelveTwo}{9.63}
\newcommand{\DwTwo}{1.78}
\newcommand{\JbTwoP}{0.079}
\newcommand{\JbTwo}{5.1}
\newcommand{\RhoOne}{0.57}
\newcommand{\RhoOneP}{$<$0.001}
\newcommand{\RhoTwo}{$-$0.08}
\newcommand{\RhoTwoP}{0.351}
\newcommand{\SimReps}{20,000}
\newcommand{\SimTrue}{1.67}
\newcommand{\SimAOls}{0.39}
\newcommand{\SimAWls}{0.57}
\newcommand{\SimACl}{94.7}
\newcommand{\SimAHc}{93.7}
\newcommand{\SimAMeanOls}{1.664}
\newcommand{\SimAMeanWls}{1.665}
\newcommand{\SimBOls}{0.68}
\newcommand{\SimBWls}{0.47}
\newcommand{\SimBCl}{72.9}
\newcommand{\SimBHc}{90.6}
\newcommand{\SimBMeanOls}{1.666}
\newcommand{\SimBMeanWls}{1.669}
\newcommand{\SimVarRatio}{43}
\newcommand{\InOlsOne}{0.22}
\newcommand{\InRfOne}{0.87}
\newcommand{\OutOlsOne}{0.15}
\newcommand{\OutRfOne}{0.00}
\newcommand{\InOlsTwo}{0.15}
\newcommand{\InRfTwo}{0.85}
\newcommand{\OutOlsTwo}{0.17}
\newcommand{\OutRfTwo}{0.01}
\newcommand{\TestFirst}{2020-08}
\newcommand{\TestLast}{2025-07}
\newcommand{\TestMonths}{60}


\begin{document}

\vspace*{14pt}
\begin{tabular}{@{}p{2.35cm}@{}p{\dimexpr\textwidth-2.35cm}@{}}
\textbf{TITLE:} & \textbf{\large When is OLS BLUE? The Gauss--Markov assumptions in a regression of Bitcoin on
macro-financial variables}\\[3pt]
\textbf{SUBJECT:} & Aim, method and results of version 2 of the \texttt{bitcoin-macro-blueml} repository, an
application of undergraduate econometrics\\[3pt]
\textbf{AUTHOR:} & Muhammed İkbal Yılmaz (myucanlar@gmail.com)
\end{tabular}

\vspace{8pt}
\textbf{Key finding.} Version 1 regressed the monthly change of the Bitcoin price in US dollars on changes of gold,
the S\&P 500, the dollar index, consumer prices and the federal funds rate, and found that OLS is not BLUE. With
the same data, \hl{the failure comes from the form of the variables, not from Bitcoin}: in dollar changes the error
variance grows with the price level (Breusch--Pagan $p = \BpOneP$, White $p = \WhiteOneP$) and the residuals are
autocorrelated (Breusch--Godfrey $p = \BgTwoOneP$); in log returns none of the checks rejects its assumption. A
random forest that explains \InRfOne{} of the variance in sample explains none of it in new months.

\section{Data and design}

\begin{itemize}
\item \textbf{Data.} The file of version 1: \Months{} months, January 2015 to July 2025, Bitcoin, gold,
the S\&P 500 and the dollar index (Investing.com), the effective federal funds rate (FRED) and the CPI-U
(FRED, BLS).\footnote{As named by version 1 (for the CPI-U, FRED in its notebook, BLS in its README); the daily value that
stands for a month was not recorded. The first date, 2015-01-01 04:50:10, is read as January 2015.}
\item \textbf{Model.} $y_t = \beta_0 + \beta_1 g_t + \beta_2 s_t + \beta_3 d_t + \beta_4 \pi_t + \beta_5 r_t + u_t$
for Bitcoin ($y$), gold ($g$), the S\&P 500 ($s$), the dollar index ($d$), consumer prices ($\pi$) and the federal
funds rate ($r$), \Nobs{} months. \textbf{Version 1:} first differences of the levels, $\Delta x_t = x_t - x_{t-1}$.
\textbf{Version 2:} log returns in per cent, $100(\ln P_t - \ln P_{t-1})$, for Bitcoin, gold, the S\&P 500 and the
dollar index; monthly inflation $100(\ln \mathrm{CPI}_t - \ln \mathrm{CPI}_{t-1})$; the change of the federal funds
rate in percentage points.
\end{itemize}

\begin{notebox}{BOX: THE GAUSS--MARKOV THEOREM}
In $y = X\beta + u$, assume (A1) linearity in $\beta$, (A2) full
column rank of $X$, (A3) $E[u \mid X] = 0$ and (A4) $\mathrm{Var}(u \mid X) = \sigma^2 I$: constant variance and no
correlation across months. Then $b = (X'X)^{-1}X'y$ has $E[b] = \beta$ and $\mathrm{Var}(b) = \sigma^2 (X'X)^{-1}$.
For any other estimator $\tilde b = Cy$ that is linear in $y$ and unbiased, $\mathrm{Var}(\tilde b) - \mathrm{Var}(b)$
is positive semi-definite: OLS is the \emph{best linear unbiased estimator}. Normality of $u$ is not needed; it makes
$t$ and $F$ exact in small samples.
\end{notebox}

\section{Analysis 1: the data and the form of the variables (Figure~\ref{fig:data})}

\textbf{Why.} The theorem concerns a linear model with constant error variance, so the form in which the variables
enter decides whether A4 can hold.

\textbf{Method.} ADF and Phillips--Perron unit-root tests (version 1's tests) on levels, differences and returns.

\begin{figure}[!htb]
\widefig{../figures/fig1_data.pdf}
\caption{(a) Bitcoin and the S\&P 500, January 2015 = 100, log scale. (b) Version 1's dependent variable, the monthly
change of Bitcoin in thousand US dollars. (c) Version 2's dependent variable, the monthly log return in per cent.}
\label{fig:data}
\end{figure}

\textbf{Finding.} Bitcoin rose from \BtcFirst{} to \BtcLast{} dollars (+\BtcGrowth\%), the S\&P 500 from \SpxFirst{} to
\SpxLast{} (+\SpxGrowth\%). ADF rejects a unit root in no level ($p \ge \LevelAdfPmin$) and in every log return ($p$
\RetAdfPmax). The levels of Bitcoin and the S\&P 500 correlate at \CorrLevels, their returns at \CorrReturns: two
trending series correlate whether or not they are related (Granger and Newbold, 1974). The change of the CPI is
ambiguous: ADF does not reject a unit root ($p = \CpiDiffAdfP$, with \InflAdfLags{} lags chosen by AIC for a series
that is not seasonally adjusted), Phillips--Perron does ($p$ \CpiDiffPpP).

\textit{Interpretation.} Panel (b) shows the problem before any test: a dollar change of Bitcoin in 2015 is measured
in tens of dollars, in 2024 in tens of thousands. Panel (c) has the same scale throughout.

\section{Analysis 2: the assumptions, one by one (Figure~\ref{fig:res}, Tables~\ref{tab:gm} and \ref{tab:coef})}

\textbf{Method.} A1: Ramsey's RESET adds $\hat y^2$ and $\hat y^3$ and tests them with $F$. A2: variance inflation
factors, $\mathrm{VIF}_j = 1/(1-R_j^2)$. A4, constant variance: Breusch--Pagan regresses $\hat u_t^2$ on the
regressors, White also on their squares and products; both use $nR^2 \sim \chi^2$. A4, no autocorrelation:
Breusch--Godfrey regresses $\hat u_t$ on the regressors and $\hat u_{t-1},\dots,\hat u_{t-p}$, $nR^2 \sim \chi^2_p$;
Durbin--Watson $\sum(\hat u_t - \hat u_{t-1})^2/\sum \hat u_t^2$ looks only at the first lag and its bounds leave
an inconclusive region. A3 has
no test: it concerns the unobserved $u$. Robust standard errors: White,
$\widehat{\mathrm{Var}}(b) = \frac{n}{n-k}(X'X)^{-1}\big(\sum_t \hat u_t^2 x_t x_t'\big)(X'X)^{-1}$, and Newey--West with
2 lags as in version 1.

\begin{table}[!htb]
\centering\small
\caption{Checks of the Gauss--Markov assumptions, \Nobs{} months. Statistic and $p$-value.}\label{tab:gm}
\begin{tabular}{llrrrr}
\toprule
 & & \multicolumn{2}{c}{Version 1: differences} & \multicolumn{2}{c}{Version 2: returns} \\
\cmidrule(lr){3-4}\cmidrule(lr){5-6}
 & Test & stat. & $p$ & stat. & $p$ \\
\midrule
A1 & RESET ($F$) & 0.57 & 0.569 & 1.04 & 0.356 \\
A2 & largest VIF & 1.73 & -- & 1.67 & -- \\
A4 & Breusch--Pagan & 11.74 & 0.039 & 2.04 & 0.844 \\
A4 & White & 36.70 & 0.013 & 13.87 & 0.837 \\
A4 & Breusch--Godfrey, 2 lags & 9.60 & 0.008 & 1.73 & 0.421 \\
A4 & Breusch--Godfrey, 6 lags & 17.14 & 0.009 & 1.76 & 0.940 \\
A4 & Durbin--Watson & 1.67 & -- & 1.78 & -- \\
-- & Jarque--Bera (normality) & 107.92 & $<$0.001 & 5.07 & 0.079 \\
\bottomrule

\end{tabular}
\end{table}

\begin{figure}[!htb]
\widefig{../figures/fig2_residuals.pdf}
\caption{Residuals of the two regressions. (a), (b) Against the Bitcoin price of the previous month (log scale).
(c), (d) Autocorrelations; grey band $\pm 1.96/\sqrt{n}$.}\label{fig:res}
\end{figure}

\textbf{Finding.} In version 1's form, A4 fails twice: the error variance depends on the regressors (Breusch--Pagan
$p = \BpOneP$, White $p = \WhiteOneP$), and the absolute residual rises with last month's Bitcoin price (rank
correlation \RhoOne, $p$ \RhoOneP); the residuals are autocorrelated (Breusch--Godfrey $p = \BgTwoOneP$ with two
lags, \BgSixOneP{} with six). In version 2's form neither part of A4 is rejected (Breusch--Pagan $p = \BpTwoP$, White
\WhiteTwoP, Breusch--Godfrey \BgTwoTwoP{} and \BgSixTwoP), nor A1 (RESET $p = \ResetTwoP$); the rank correlation with
the price is \RhoTwo{} ($p = \RhoTwoP$). Collinearity is no issue in either form (largest VIF \VifOne{} and \VifTwo).
Normality is rejected for version 1 (Jarque--Bera $p$ \JbOneP) but not for version 2 ($p = \JbTwoP$).

\begin{table}[!htb]
\centering\small
\caption{Version 2: coefficients with classical, White (HC1) and Newey--West (2 lags) standard errors; $p$-value with
White standard errors. $R^2 = \RsqTwo$.}\label{tab:coef}
\setlength{\tabcolsep}{4pt}\begin{tabular}{lrrrrr}
\toprule
 & coef. & classical & White & Newey--West & $p$ \\
\midrule
constant & 4.77 & 2.12 & 2.24 & 2.48 & 0.033 \\
gold & 0.49 & 0.50 & 0.45 & 0.40 & 0.280 \\
S\&P 500 & 1.67 & 0.39 & 0.36 & 0.33 & $<$0.001 \\
dollar index & 0.93 & 1.09 & 0.96 & 0.93 & 0.331 \\
inflation & $-$5.68 & 5.12 & 5.36 & 5.80 & 0.289 \\
fed funds rate (change) & $-$7.94 & 9.47 & 8.98 & 9.13 & 0.377 \\
\bottomrule

\end{tabular}
\end{table}

\textit{Interpretation.} With A1, A2 and A4 not rejected, OLS on returns is BLUE provided A3 holds. Only the S\&P 500
matters: a month in which its log return is 1 point higher goes with a Bitcoin log return \SpxTwo{} points higher (White s.e.\ \SpxTwoSeHc);
the other variables have $p \ge \OthersTwoPmin$. The three standard errors differ little, as expected when A4 holds.

\section{Analysis 3: what BLUE means, in a simulation (Figure~\ref{fig:sim})}

\textbf{Method.} The regressors of version 2 are held fixed, $\beta$ and $\sigma$ are set to their estimates, and
\SimReps{} samples $y = X\beta + u$ are drawn under (A) $u_t \sim N(0,\sigma^2)$ and (B)
$\mathrm{Var}(u_t) \propto 0.25 + z_t^2$, $z_t$ the standardised S\&P 500 return, with the same average variance.
Two linear unbiased estimators of $\beta_2$ are compared: OLS and weighted least squares,
$b_W = (X'WX)^{-1}X'Wy$ with $W = \mathrm{diag}\{1/(0.25 + z_t^2)\}$, which is generalised least squares under B.

\begin{figure}[!htb]
\widefig{../figures/fig3_simulation.pdf}
\caption{(a), (b) Distribution of the OLS and WLS estimates of the S\&P 500 coefficient under A and B; red: true
value \SimTrue. (c) Share of samples in which the 95\% interval of OLS contains the true value.}\label{fig:sim}
\end{figure}

\textbf{Finding.} Both estimators are unbiased under both laws (means \SimAMeanOls{} and \SimAMeanWls{} under A,
\SimBMeanOls{} and \SimBMeanWls{} under B). Under A, OLS has the smaller standard deviation (\SimAOls{} against
\SimAWls), as the theorem says. Under B, WLS does (\SimBWls{} against \SimBOls), and the classical 95\% interval of OLS
contains the true value in only \SimBCl\% of samples; White's interval in \SimBHc\%.

\textit{Interpretation.} A failure of A4 does not bias OLS; it makes it inefficient and its classical standard
errors wrong. White's standard errors repair most, not all, of the coverage in a sample of this size.

\section{Analysis 4: the random forest (Figure~\ref{fig:rf})}

\textbf{Method.} Version 1's forest (500 trees, fully grown, seed 42) and OLS, first on all months, then for each of
the last \TestMonths{} months (\TestFirst{} to \TestLast) fitted on the earlier months only and applied to that
month: $R^2_{\mathrm{new}} = 1 - \sum(y_t - \hat y_t)^2/\sum(y_t - \bar y_t)^2$, $\bar y_t$ the mean of the earlier
months.

\begin{figure}[!htb]
\widefig{../figures/fig4_forest.pdf}
\caption{$R^2$ of OLS and of the random forest in sample and on months not used in fitting.}\label{fig:rf}
\end{figure}

\textbf{Finding.} In sample the forest reaches $R^2 = \InRfOne$ (version 1's form) and \InRfTwo{} (returns), against
\InOlsOne{} and \InOlsTwo{} for OLS. On new months the forest falls to \OutRfOne{} and \OutRfTwo, OLS keeps \OutOlsOne{}
and \OutOlsTwo.

\textit{Interpretation.} With \Nobs{} months, fully grown trees memorise the sample; the in-sample $R^2$ of version 1
measured that. The linear model carries over to new months; the forest finds no additional structure.

\section{Conclusions and limitations}

\begin{enumerate}
\item Whether OLS is BLUE depends on the form of the variables: in dollar changes A4 fails, in log returns it is not
rejected.
\item Only the S\&P 500 is related to Bitcoin in the same month; gold, the dollar, inflation and the policy rate are
not.
\item A3 cannot be tested. Bitcoin and the S\&P 500 react to the same news in the same month, so the coefficient
describes co-movement, not the effect of one on the other; version 1's Granger tests find no lead of the S\&P 500
over Bitcoin ($p \ge \GrangerPmin$ for 1--6 lags).
\end{enumerate}

\textbf{Scope and verification note.} Every number is written from \texttt{figures/*.csv} by \texttt{make\_paper.py};
the unit tests check the computations against matrix formulas, auxiliary regressions written out, exact variances
for the simulation and the printed output of version 1, whose code reproduces it on the stored file. The sample is
\Nobs{} months; the monthly timing of the source series was not recorded.

\section*{References}
{\small
Breiman, L. (2001). Random forests. \textit{Machine Learning}, 45(1), 5--32.\\
Breusch, T. S. (1978). Testing for autocorrelation in dynamic linear models. \textit{Australian Economic Papers},
17(31), 334--355.\\
Breusch, T. S., Pagan, A. R. (1979). A simple test for heteroscedasticity and random coefficient variation.
\textit{Econometrica}, 47(5), 1287--1294.\\
Godfrey, L. G. (1978). Testing against general autoregressive and moving average error models when the regressors
include lagged dependent variables. \textit{Econometrica}, 46(6), 1293--1301.\\
Granger, C. W. J., Newbold, P. (1974). Spurious regressions in econometrics. \textit{Journal of Econometrics}, 2(2),
111--120.\\
Newey, W. K., West, K. D. (1987). A simple, positive semi-definite, heteroskedasticity and autocorrelation consistent
covariance matrix. \textit{Econometrica}, 55(3), 703--708.\\
Ramsey, J. B. (1969). Tests for specification errors in classical linear least-squares regression analysis.
\textit{Journal of the Royal Statistical Society, Series B}, 31(2), 350--371.\\
White, H. (1980). A heteroskedasticity-consistent covariance matrix estimator and a direct test for
heteroskedasticity. \textit{Econometrica}, 48(4), 817--838.\\
Wooldridge, J. M. (2020). \textit{Introductory Econometrics: A Modern Approach} (7th ed.). Cengage.

\smallskip
An AI coding assistant was used in writing the code and in preparing this note.
}

\end{document}
